\section{Kliegel--Levine Axisymmetric Transonic Initial-Line Model}
\label{sec:kliegel_levine_initial_line}
% =========================================================

\subsection{Purpose of the model}

The Kliegel--Levine model supplies a local analytical approximation to the
axisymmetric flow field close to the nozzle throat.  It is not a marching
method and it does not solve a computational mesh.  Instead, it provides the
continuous functions

\begin{equation}
    U(X,Y)=\frac{V_x}{a^*},
    \qquad
    V(X,Y)=\frac{V_r}{a^*},
\end{equation}

which can subsequently be sampled on a wholly supersonic initial-data line.
The state stored at each selected point is

\begin{equation}
    \boldsymbol{W}_j=
    \left[
        x_j,\ r_j,\ M_j,\ \theta_j,\ p_j,\ T_j,\ \rho_j,
        \ V_{x,j},\ V_{r,j}
    \right]^T.
\end{equation}

The approximation assumes steady, inviscid, irrotational and isentropic flow
of a calorically perfect gas with locally constant \(\gamma\).  The analytical
series is only intended for the transonic throat neighbourhood.

\subsection{Local geometry and nondimensional coordinates}

The throat radius and the convergent-side radius of curvature define

\begin{equation}
    R_t,
    \qquad
    R_{c,\mathrm{sub}},
    \qquad
    R_\kappa=\frac{R_{c,\mathrm{sub}}}{R_t}.
\end{equation}

For the present nozzle,

\begin{equation}
    R_{c,\mathrm{sub}}=1.5R_t,
    \qquad
    R_\kappa=1.5.
\end{equation}

The divergent-side throat arc remains part of the prescribed nozzle contour
but does not define the Kliegel--Levine expansion parameter.  The local
coordinates are

\begin{equation}
    X=\frac{x-x_t}{R_t},
    \qquad
    Y=\frac{r}{R_t},
\end{equation}

and the scaled axial coordinate is

\begin{equation}
    \boxed{
    Z=X\sqrt{\frac{2R_\kappa}{\gamma+1}}
    }.
\end{equation}

For compactness, define

\begin{equation}
    \Lambda=1+R_\kappa,
    \qquad
    \varepsilon_\kappa=\frac{1}{\Lambda}.
\end{equation}

The symbol \(R_g\) is reserved for the specific gas constant and is therefore
not confused with the curvature ratio \(R_\kappa\).

The geometric visualization of this initial transient line it's represented on figure \ref{fig:kl_actual_geometry_coordinates}.

\clearpage
\begin{landscape}

\begin{figure}[p]
\centering

\def\KLGeometryCSV{geometry.csv}
\def\KLInitialCSV{moc_initial_data_line.csv}

\resizebox{0.995\linewidth}{!}{
\begin{tikzpicture}
\begin{axis}[
    width=26.8cm,
    height=16.0cm,
    xmin=-2,
    xmax=121,
    ymin=-52,
    ymax=52,
    axis lines=middle,
    xlabel={global axial coordinate, \(x\) [mm]},
    ylabel={radial coordinate, \(r\) [mm]},
    tick label style={font=\small},
    label style={font=\normalsize},
    grid=none,
    minor grid style={gray!15},
    major grid style={gray!25},legend style={
    at={(0.23,0.88)},
    anchor=south,
    legend columns=1,
    draw=none,
    fill=none,
    font=\small
},
legend cell align={left},
    clip=false
]

% Actual upper and lower nozzle contours from geometry.csv.
\addplot[black,very thick]
table[x=x_mm,y=radius_mm,col sep=comma] {\KLGeometryCSV};

\addplot[black,very thick,forget plot]
table[x=x_mm,y expr=-\thisrow{radius_mm},col sep=comma] {\KLGeometryCSV};

% Actual sampled initial-data line from moc_initial_data_line.csv.
\addplot[red!75!black,very thick]
table[
    x expr=1000*\thisrow{x_m},
    y expr=1000*\thisrow{radius_m},
    col sep=comma
] {\KLInitialCSV};


% Geometric throat and local origin used by X=(x-x_t)/R_t.
\addplot[gray!75,densely dashed,thick]
coordinates {(46.0447788,-48) (46.0447788,48)};

\addplot[only marks,mark=*,mark size=2.4pt,black,forget plot]
coordinates {(46.0447788,17.28) (46.0447788,-17.28)};

\fill (axis cs:46.0447788,0) circle (2.1pt);
\node[anchor=north east,fill=white,inner sep=2pt]
at (axis cs:45.3,-5.0)
{local origin \(O_t=(x_t,0)\)};

%\node[anchor=south east,fill=white,inner sep=2pt,align=right]
%at (axis cs:45.2,18.5)
%{throat wall point \(T\)\\
%\(x_t=46.0448\,\mathrm{mm}\), \(R_t=17.28\,\mathrm{mm}\)};

% Throat-radius construction.
\draw[<->,thick]
    (axis cs:42.3,0) -- (axis cs:42.3,17.28)
    node[midway,left,fill=white,inner sep=1.5pt]
    {\(R_t\)};
\draw[thin]
    (axis cs:42.3,17.28) -- (axis cs:46.0447788,17.28);

% Convergent-side curvature centre and radius.
\fill (axis cs:46.0447788,43.2) circle (1.8pt);
\draw[densely dashed,thick,orange!80!black]
    (axis cs:46.0447788,43.2) -- (axis cs:46.0447788,17.28);
\node[anchor=west,fill=white,inner sep=2pt,text=orange!70!black,align=left]
at (axis cs:47.4,35.2)
{\(R_{c,\mathrm{sub}}=1.5R_t\)\\
local convergent curvature};
\node[anchor=south east,fill=white,inner sep=1.5pt]
at (axis cs:45.4,43.2)
{\(C_{\mathrm{sub}}\)};

% Arbitrary point P in the local analytical field.  This is a geometrical
% construction, not a plotted solver result.
\coordinate (KLP) at (axis cs:50.3176,9.0);
\coordinate (KLPx) at (axis cs:50.3176,14.0);
\coordinate (KLPy) at (axis cs:46.0447788,9.0);
\fill[blue!75!black] (KLP) circle (2.2pt);
\node[anchor=south west,text=blue!75!black,fill=white,inner sep=2pt]
at (KLP) {arbitrary local point \(P(x,r)\)};

\draw[blue!60!black,densely dashed,thick] (KLP) -- (KLPx);
\draw[blue!60!black,densely dashed,thick] (KLP) -- (KLPy);

% Horizontal and radial dimensions defining X and Y.
\draw[<->,very thick,blue!75!black]
    (axis cs:46.0447788,12.0) -- (axis cs:50.3176,12.0)
    node[midway,above=3.5pt,fill=white,inner sep=1.5pt]
    {\(x-x_t=XR_t\)};

\draw[<->,very thick,blue!75!black]
    (axis cs:50.3176,0.0) -- (axis cs:50.3176,9.0)
    node[midway,right=2pt,fill=white,inner sep=1.5pt]
    {\(r=YR_t\)};

% Compact definition boxes.
\node[
    draw,
    rounded corners,
    thick,
    fill=white,
    align=center,
    inner sep=5pt,
    anchor=north west
] at (axis cs:64,59)
{
    \textbf{Local Kliegel--Levine coordinates}\\[4pt]
    \(\displaystyle
    X=\frac{x-x_t}{R_t},
    \qquad
    Y=\frac{r}{R_t}
    \)\\[4pt]
    \(\displaystyle
    Z=X\sqrt{\frac{2R_\kappa}{\gamma+1}}
    \)
};

\node[
    draw,
    rounded corners,
    thick,
    fill=white,
    align=left,
    inner sep=5pt,
    anchor=south west
] at (axis cs:82,-49)
{\textbf{Curvature parameter}\\[4pt]
\(\displaystyle
R_\kappa=\frac{R_{c,\mathrm{sub}}}{R_t},\qquad
\Lambda=1+R_\kappa,\qquad
\varepsilon_\kappa=\frac{1}{\Lambda}
\)};

\draw[-{Latex[length=3mm]},very thick]
    (axis cs:91,21) -- (axis cs:108,21)
    node[midway,above] {flow direction};
\addlegendimage{line legend,red!75!black,very thick}
\addlegendentry{Kliegel--Levine initial-data line}

\addlegendimage{line legend,black,very thick}
\addlegendentry{Nozzle contour}
\end{axis}
\end{tikzpicture}
}

\caption{Geometrical definition of the Kliegel--Levine local coordinates on
the actual nozzle contour read from \texttt{geometry.csv}.  The point \(P\) is
an arbitrary location in the local analytical field and is included only to
illustrate \(X=(x-x_t)/R_t\) and \(Y=r/R_t\); it is not a computed initial-line
result.  The vertical scale is compressed to obtain a clearer horizontal
layout.}
\label{fig:kl_actual_geometry_coordinates}

\end{figure}
\end{landscape}


\subsection{Meaning of the series coefficients}

Before their values are known, the first-order polynomials may be represented
as

\begin{equation}
    u_1=A_0+A_2Y^2+A_ZZ,
\end{equation}

\begin{equation}
    v_1=Y\left(B_0+B_2Y^2+B_ZZ\right).
\end{equation}

The symbols \(A_0,A_2,A_Z,B_0,B_2,B_Z\) are temporary unknown weights of
the polynomial basis.  They are not empirical efficiency factors and they are
not supplied as simulation inputs.  In the analytical derivation, the assumed
polynomials are inserted into the transonic equations and the coefficients of
equal powers of \(Y\) and \(Z\) are matched.  Axis symmetry and throat-wall
tangency complete the resulting algebraic system.

With the Kliegel--Levine normalization used here, the solution of that
coefficient system is

\begin{equation}
    A_0=-\frac14,
    \qquad
    A_2=\frac12,
    \qquad
    A_Z=1,
\end{equation}

\begin{equation}
    B_0=-\frac14,
    \qquad
    B_2=\frac14,
    \qquad
    B_Z=1.
\end{equation}

Consequently,

\begin{equation}
    \boxed{
    u_1=\frac12Y^2-\frac14+Z
    },
\end{equation}

\begin{equation}
    \boxed{
    v_1=\frac14Y^3-\frac14Y+YZ
    }.
\end{equation}

These numerical fractions are general for this first-order axisymmetric
Kliegel--Levine normalization; they are not specific to
\(R_{c,\mathrm{sub}}=1.5R_t\).  A change in throat curvature changes
\(\Lambda\), \(\varepsilon_\kappa\) and \(Z\), while the displayed
first-order fractions remain unchanged.  A different coordinate scaling,
planar geometry or expansion family would produce a different coefficient
set.  Higher-order coefficients additionally depend on \(\gamma\).

The polynomial parity incorporates the axis conditions directly.  The axial
polynomials contain only even powers of \(Y\), whereas the radial polynomials
contain only odd powers.  Therefore

\begin{equation}
    \left.\frac{\partial U}{\partial Y}\right|_{Y=0}=0,
    \qquad
    V(X,0)=0.
\end{equation}

At the geometric throat wall, \(Y=1\) and \(Z=0\), and the first-order radial
polynomial gives

\begin{equation}
    v_1(0,1)=\frac14-\frac14=0,
\end{equation}

consistent with the zero local wall slope at the throat.  These conditions are
embedded in the closed-form coefficients; the implementation does not impose
them again through a separate numerical boundary-condition solve.

\subsection{Third-order velocity expansion}

The second-order polynomials are

\begin{align}
u_2={}&
    \frac{2\gamma+9}{24}Y^4
    -\frac{4\gamma+15}{24}Y^2
    +\frac{10\gamma+57}{288}
    +Z\left(Y^2-\frac58\right)
    -\frac{2\gamma-3}{6}Z^2,
\end{align}

\begin{align}
v_2={}&
    \frac{\gamma+3}{9}Y^5
    -\frac{20\gamma+63}{96}Y^3
    +\frac{28\gamma+93}{288}Y \nonumber\\
    &+Z\left[
        \frac{2\gamma+9}{6}Y^3
        -\frac{4\gamma+15}{12}Y
    \right]
    +YZ^2.
\end{align}

The third-order polynomials used by the implementation are

\begin{align}
u_3={}&
 \frac{556\gamma^2+1737\gamma+3069}{10368}Y^6
-\frac{388\gamma^2+1161\gamma+1881}{2304}Y^4 \nonumber\\
&+\frac{304\gamma^2+831\gamma+1242}{1728}Y^2
-\frac{2708\gamma^2+7839\gamma+14211}{82944} \nonumber\\
&+Z\left[
 \frac{52\gamma^2+51\gamma+327}{384}Y^4
-\frac{52\gamma^2+75\gamma+279}{192}Y^2
+\frac{92\gamma^2+180\gamma+639}{1152}
\right] \nonumber\\
&+Z^2\left[-\frac{7\gamma-3}{8}Y^2
+\frac{13\gamma-27}{48}\right]
+\frac{4\gamma^2-57\gamma+27}{144}Z^3,
\end{align}

\begin{align}
v_3={}&
 \frac{6836\gamma^2+23031\gamma+30627}{82944}Y^7
-\frac{3380\gamma^2+11391\gamma+15291}{13824}Y^5 \nonumber\\
&+\frac{3424\gamma^2+11271\gamma+15228}{13824}Y^3
-\frac{7100\gamma^2+22311\gamma+30249}{82944}Y \nonumber\\
&+Z\left[
 \frac{556\gamma^2+1737\gamma+3069}{1728}Y^5
-\frac{388\gamma^2+1161\gamma+1181}{576}Y^3
+\frac{304\gamma^2+831\gamma+1242}{864}Y
\right] \nonumber\\
&+Z^2\left[
 \frac{52\gamma^2+51\gamma+327}{192}Y^3
-\frac{52\gamma^2+75\gamma+279}{192}Y
\right]
-\frac{7\gamma-3}{12}YZ^3.
\end{align}

The nondimensional velocity components are assembled as

\begin{equation}
\boxed{
U=
1+\frac{u_1}{\Lambda}
+\frac{u_1+u_2}{\Lambda^2}
+\frac{u_1+2u_2+u_3}{\Lambda^3}
},
\label{eq:kl_axial_velocity}
\end{equation}

\begin{equation}
\boxed{
V=
\sqrt{\frac{\gamma+1}{2\Lambda}}
\left[
\frac{v_1}{\Lambda}
+\frac{\frac32v_1+v_2}{\Lambda^2}
+\frac{\frac{15}{8}v_1+\frac52v_2+v_3}{\Lambda^3}
\right]
}.
\label{eq:kl_radial_velocity}
\end{equation}

No coefficient system is solved during program execution.  Equations
\eqref{eq:kl_axial_velocity} and \eqref{eq:kl_radial_velocity} are evaluated
directly at each requested \((X,Y)\) location.

\subsection{Physical state recovered from the analytical field}

The critical velocity is

\begin{equation}
    a^*=\sqrt{\frac{2\gamma R_gT_0}{\gamma+1}}.
\end{equation}

After evaluating \(U\) and \(V\), define

\begin{equation}
    \overline q=\sqrt{U^2+V^2}.
\end{equation}

The local sound-speed ratio and Mach number are

\begin{equation}
    \frac{a^2}{a^{*2}}
    =\frac12\left[\gamma+1-(\gamma-1)\overline q^2\right],
\end{equation}

\begin{equation}
    \boxed{
    M=\frac{\overline q}
    {\sqrt{\frac12[\gamma+1-(\gamma-1)\overline q^2]}}
    }.
\end{equation}

The flow angle and dimensional velocity components are

\begin{equation}
    \theta=\operatorname{atan2}(V,U),
    \qquad
    V_x=a^*U,
    \qquad
    V_r=a^*V.
\end{equation}

The static properties follow from

\begin{equation}
    T=\frac{T_0}{1+\frac12(\gamma-1)M^2},
\end{equation}

\begin{equation}
    p=p_0\left(\frac{T}{T_0}\right)^{\gamma/(\gamma-1)},
    \qquad
    \rho=\frac{p}{R_gT}.
\end{equation}

\subsection{Sampling the curved initial-data line}

The analytical field itself is continuous and does not require a mesh.  The
integer \(N_r\) only determines how many samples are retained on the initial
line.  Points are first ordered from the throat wall to the symmetry axis by

\begin{equation}
    Y_j=
    \sin^{3/2}\left[
        \frac{\pi}{2}\frac{N_r-1-j}{N_r-1}
    \right],
    \qquad j=0,\ldots,N_r-1.
\end{equation}

The first point is

\begin{equation}
    X_0=0,
    \qquad
    Y_0=1.
\end{equation}

For each subsequent radial sample, the curve is continued using the local
minus-family direction evaluated at the preceding point:

\begin{equation}
    \mu_{j-1}=\sin^{-1}\left(\frac{1}{M_{j-1}}\right),
\end{equation}

\begin{equation}
    \boxed{
    X_j=X_{j-1}
    +\frac{Y_j-Y_{j-1}}
    {\tan(\theta_{j-1}-\mu_{j-1})}
    }.
\end{equation}

The Kliegel--Levine polynomials are evaluated at every newly obtained
\((X_j,Y_j)\).  This recurrence selects the curve on which the analytical
field is sampled; it does not solve the Kliegel--Levine field itself.  Physical
coordinates are finally recovered from

\begin{equation}
    x_j=x_t+R_tX_j,
    \qquad
    r_j=R_tY_j.
\end{equation}


\subsection{Applicability and implementation note}

The coefficient values are fixed by the selected Kliegel--Levine formulation
and are stored directly in the implementation.  The case-specific quantities
are \(R_t\), \(R_{c,\mathrm{sub}}/R_t\), \(\gamma\), \(R_g\), \(p_0\) and
\(T_0\).  The model does not reproduce viscous effects, shocks, variable
specific heats or the complete divergent flow.  Its validity should therefore
be assessed through positivity, a wholly supersonic sampled line and the
subsequent mass-flow consistency checks.

% ---------------------------------------------------------
% FIGURE 1: ACTUAL CSV GEOMETRY AND KLIEGEL--LEVINE LINE
% ---------------------------------------------------------

% ---------------------------------------------------------
% FIGURE 1: ACTUAL CSV GEOMETRY AND KL COORDINATE CONSTRUCTION
% ---------------------------------------------------------




% ---------------------------------------------------------
% FIGURE 2: HORIZONTAL KLIEGEL--LEVINE ALGORITHM
% ---------------------------------------------------------

\clearpage
\begin{landscape}


\subsection{Kiegle-Levine Initial Transient Line Model Workflow}
\begin{figure}[H]
\centering

\resizebox{0.985\linewidth}{!}{
\begin{tikzpicture}[
    >=Latex,
    node distance=0.75cm and 0.65cm,
    process/.style={
        rectangle,
        draw,
        rounded corners,
        thick,
        align=center,
        text width=4.55cm,
        minimum height=2.45cm,
        inner sep=5pt,
        font=\small
    },
    wideprocess/.style={
        rectangle,
        draw,
        rounded corners,
        thick,
        align=center,
        text width=6.0cm,
        minimum height=2.45cm,
        inner sep=5pt,
        font=\small
    },
    startstop/.style={
        rectangle,
        draw,
        double,
        rounded corners,
        very thick,
        align=center,
        text width=4.8cm,
        minimum height=2.45cm,
        inner sep=6pt,
        font=\small
    },
    decision/.style={
        diamond,
        draw,
        thick,
        aspect=2.4,
        align=center,
        text width=3.0cm,
        inner sep=2pt,
        font=\small
    },
    flow/.style={->,very thick,rounded corners}
]

% First horizontal row.
\node[startstop] (inputs)
{\textbf{Physical and numerical inputs}
\[
\gamma,\ R_g,\ p_0,\ T_0,
\ R_t,\ R_{c,\mathrm{sub}},\ N_r
\]};

\node[process,below=of inputs] (scales)
{\textbf{1. Form the K--L scales}
\[
R_\kappa=\frac{R_c}{R_t},\quad
\Lambda=1+R_\kappa
\]
\[
a^*=\sqrt{\frac{2\gamma R_gT_0}{\gamma+1}}
\]};

\node[wideprocess,right=of scales] (sampling)
{\textbf{2. Select radial samples}
\[
Y_j=\sin^{3/2}\!\left[
\frac{\pi}{2}\frac{N_r-1-j}{N_r-1}
\right]
\]
Start at \(X_0=0,\ Y_0=1\).};

\node[wideprocess,right=of sampling] (curve)
{\textbf{3. Locate the next point}
\[
\mu_{j-1}=\sin^{-1}(1/M_{j-1})
\]
\[
X_j=X_{j-1}+\frac{Y_j-Y_{j-1}}
{\tan(\theta_{j-1}-\mu_{j-1})}
\]};

% Second horizontal row, traversed from right to left.
\node[wideprocess,below=1.25cm of curve] (polynomials)
{\textbf{4. Evaluate the third-order polynomials}
\[
Z_j=X_j\sqrt{\frac{2R_\kappa}{\gamma+1}}
\]
Compute \(u_1,u_2,u_3\) and
\(v_1,v_2,v_3\) at \((Z_j,Y_j)\).};

\node[wideprocess,left=of polynomials] (velocity)
{\textbf{5. Assemble the analytical velocity}
\[
U_j=1+\frac{u_1}{\Lambda}
+\frac{u_1+u_2}{\Lambda^2}
+\frac{u_1+2u_2+u_3}{\Lambda^3}
\]
Evaluate the corresponding expression for \(V_j\).};

\node[wideprocess,left=of velocity] (flowstate)
{\textbf{6. Recover Mach number and angle}
\[
\bar q_j=\sqrt{U_j^2+V_j^2}
\]
\[
M_j=\frac{\bar q_j}
{\sqrt{[\gamma+1-(\gamma-1)\bar q_j^2]/2}}
\]
\[
\theta_j=\operatorname{atan2}(V_j,U_j)
\]}

\node[decision,left=of flowstate] (valid)
{Is the point physical?
\[
M_j>1,\qquad a_j^2>0
\]};

% Third horizontal row, traversed from left to right.
\node[wideprocess,below=1.25cm of valid] (thermo)
{\textbf{7. Recover dimensional state}
\[
V_{x,j}=a^*U_j,\quad V_{r,j}=a^*V_j
\]
\[
T_j,\quad p_j,\quad \rho_j=\frac{p_j}{R_gT_j}
\]};

\node[process,right=of thermo] (store)
{\textbf{8. Store point \(j\)}
\[
\boldsymbol W_j=
[x_j,r_j,M_j,\theta_j,p_j,T_j,
\rho_j,V_{x,j},V_{r,j}]^T
\]};

\node[decision,right=of store] (finished)
{Final radial sample?
\[
j=N_r-1
\]};

\node[startstop,right=of finished] (output)
{\textbf{Complete Kliegel--Levine initial line}
\[
\mathcal L_i=
\{\boldsymbol W_0,\ldots,
\boldsymbol W_{N_r-1}\}
\]};

% Main snake-shaped flow.
\draw[flow] (inputs) -- (scales);
\draw[flow] (scales) -- (sampling);
\draw[flow] (sampling) -- (curve);
\draw[flow] (curve) -- (polynomials);
\draw[flow] (polynomials) -- (velocity);
\draw[flow] (velocity) -- (flowstate);
\draw[flow] (flowstate) -- (valid);
\draw[flow] (valid) -- node[right] {yes} (thermo);
\draw[flow] (thermo) -- (store);
\draw[flow] (store) -- (finished);
\draw[flow] (finished) -- node[above] {yes} (output);

% Loop to the next sampled point.
\draw[flow]
    (finished.north)
    -- ++(0,0.5) -- node[pos=0.25,above] {no: \(j\leftarrow j+1\)}++(14,0) -- ++(0, 6.6)
    -- 
    (curve.east);

% Invalid-state branch.
\node[process,above=1.0cm of valid] (invalid)
{Reject the sampled line or revise the local transonic approximation. There could be points in a line that fulfill the requisites. However, if one points does not meet the requirements the line should be discarded and another line should be chosen.};
\draw[flow]
    (valid.north)--
    (invalid.south);

\end{tikzpicture}
}

\caption{Horizontal computational procedure used to sample the third-order
Kliegel--Levine analytical field on its curved supersonic initial-data line.
The radial stations sample a continuous analytical solution; they do not form
a mesh on which the Kliegel--Levine equations are numerically solved.}
\label{fig:kl_initial_line_algorithm}

\end{figure}
\end{landscape}

